% ===========================================================================
% Section 1 -- Problem description and mathematical models
% ===========================================================================

\subsection{Notation}
\label{sec:notation}

\paragraph{Duties.}
Let $\D=\{1,\dots,n\}$ be the set of driving duties. Duty $i$ starts at time
$s_i$ and ends at time $e_i>s_i$; its driving time is $d_i=e_i-s_i$. For every
ordered pair $(i,j)$ of distinct duties, $t_{ij}\ge 0$ is the time needed to
move from the place where $i$ ends to the place where $j$ starts. All times are
integers (minutes) in the horizon $[0,H]$, with $H=\max_i e_i$.

\paragraph{Compatibility graph.}
Duty $j$ can be performed after duty $i$ by the same employee if and only if
the employee reaches the start of $j$ in time:
\begin{equation}\label{eq:arc}
  (i,j)\in\A \iff e_i+t_{ij}\le s_j .
\end{equation}
$G=(\D,\A)$ is the (given) compatibility graph. Since $(i,j)\in\A$ implies
$s_i<e_i\le s_j$, sorting duties by start time gives a topological order, so
$G$ is a \emph{directed acyclic graph} (DAG).

\paragraph{Employee types.}
There are two types $\tau\in\{\mathrm{F},\mathrm{H}\}$ (full-time,
half-time), each in unlimited supply, with
\[
  (W_{\mathrm F},T_{\mathrm F},c_{\mathrm F})=(W,\;T,\;1),\qquad
  (W_{\mathrm H},T_{\mathrm H},c_{\mathrm H})=(W/2,\;T/2,\;0.6),
\]
where $W_\tau$ is the maximum working time, $T_\tau$ the maximum total time
and $c_\tau$ the cost of one employee of type $\tau$ (normalised so that a
full-time employee costs $1$).

\paragraph{Shifts.}
A \emph{shift} is the ordered sequence of duties $p=(i_1,\dots,i_m)$
performed by one employee. It is \emph{feasible for type $\tau$} if
\begin{enumerate}[label=(S\arabic*)]
  \item $p$ is a path in $G$: $(i_\ell,i_{\ell+1})\in\A$ for
        $\ell=1,\dots,m-1$;
  \item working time: $\sum_{\ell=1}^m d_{i_\ell}\le W_\tau$;
  \item total time (``spread''): $e_{i_m}-s_{i_1}\le T_\tau$.
\end{enumerate}
We write $\Omega_\tau$ for the set of feasible shifts of type $\tau$. The
problem is to partition $\D$ into shifts, give each shift a type for which it
is feasible, and minimise the total cost $\sum c_\tau$ over the shifts used.

\subsection{Assumptions and interpretation choices}
\label{sec:assumptions}

The assignment leaves some details open. We fix them as follows (items marked
($\ast$) are ones we would like the lecturer to confirm).

\begin{enumerate}[label=(A\arabic*)]
  \item\label{ass:work} ($\ast$) \emph{Working time counts only the driving
    of the duties}, $\sum_{i\in p}d_i$. Deadhead movements $t_{ij}$ between
    duties are not counted, although they do count in the total time, since
    they lie between the start of the first duty and the end of the last.
    If deadheads must count too, formulation~(AF) below handles it by adding
    $t_{ij}$ to the arc weights in~\eqref{af:work}.
  \item\label{ass:cost} \emph{Cost is a fixed amount per employee used}
    ($1$ or $0.6$), independent of how many hours they actually work.
  \item\label{ass:wt} \emph{$W<T$.} Working time never exceeds total time
    (duties do not overlap), so if $W\ge T$ the working-time limit is redundant.
    The interesting case is $W<T$, and all our instances satisfy it.
  \item\label{ass:single} \emph{Every duty can be covered alone by a full-time
    employee}, i.e.\ $d_i\le W$ for all $i$; otherwise the instance is
    infeasible. Duties with $d_i>W/2$ can only be covered by full-time
    employees.
  \item\label{ass:tri} \emph{Triangle inequality (only where stated).}
    Formulation~(AF) works for an arbitrary DAG $G$, as the assignment speaks
    of a ``given'' graph. Formulation~(CF) additionally needs $G$ to be
    transitive. This is guaranteed when
    $t_{ij}\le t_{ik}+t_{kj}$, which holds whenever transfer times are
    shortest-path or Euclidean travel times between locations
    (Proposition~\ref{prop:transitive}).
\end{enumerate}

\begin{proposition}\label{prop:transitive}
If the transfer times satisfy $t_{ik}\le t_{ij}+t_{jk}$ for all $i,j,k$, then
$G$ is transitive.
\end{proposition}
\begin{proof}
Let $(i,j),(j,k)\in\A$. Then
$e_i+t_{ik}\le e_i+t_{ij}+t_{jk}\le s_j+t_{jk}\le e_j+t_{jk}\le s_k$,
hence $(i,k)\in\A$.
\end{proof}

\begin{remark}
Every half-time shift is also a feasible full-time shift
($\Omega_{\mathrm H}\subseteq\Omega_{\mathrm F}$), so a shift in
$\Omega_{\mathrm H}$ is always given to the cheaper half-time employee. Per
unit of working time, full-time employees are cheaper ($1/W$ against
$1.2/W$). Half-time employees pay off only for ``short'' shifts, those with
working time at most $W/2$ and spread at most $T/2$, that a full-time employee
would leave largely idle. This trade-off is what makes the mix of types
non-trivial.
\end{remark}

\subsection{Compact formulations}
\label{sec:compact}

Both compact models use an explicit pool of candidate employees
$\K=\K_{\mathrm F}\cup\K_{\mathrm H}$, with $|\K_\tau|=U_\tau$. Employee
$k\in\K_\tau$ has $(W_k,T_k,c_k)=(W_\tau,T_\tau,c_\tau)$. The pool must be
large enough not to cut off every optimal solution. $U_{\mathrm F}=n$ is
always valid, and so is $U_{\mathrm H}=|\{i:d_i\le W/2\}|$. Given the cost
$\bar z$ of any feasible solution (e.g.\ from a greedy heuristic), an optimal
solution uses at most $\lfloor\bar z\rfloor$ full-time and
$\lfloor\bar z/0.6\rfloor$ half-time employees, which gives much smaller
pools.

\subsubsection{Arc-flow formulation (AF)}

Add a source $\sigma$ and a sink $\omega$ to $G$:
$\bar\A=\A\cup\{(\sigma,i):i\in\D\}\cup\{(i,\omega):i\in\D\}$. For a node
$v$, $\delta^-(v)$ and $\delta^+(v)$ are its entering and leaving arcs in
$\bar\A$. The binary variable $z^k_a$ equals $1$ if employee $k$ uses
arc $a\in\bar\A$. Arc $(\sigma,i)$ means ``$i$ is the first duty of $k$'',
arc $(i,\omega)$ means ``$i$ is the last one''.
\begin{align}
\text{(AF)}\quad
\min\;& \sum_{k\in\K} c_k \sum_{i\in\D} z^k_{\sigma i} \label{af:obj}\\
\text{s.t.}\;
& \sum_{k\in\K}\sum_{a\in\delta^-(i)} z^k_a = 1
      && i\in\D \label{af:cover}\\
& \sum_{a\in\delta^-(i)} z^k_a-\sum_{a\in\delta^+(i)} z^k_a = 0
      && i\in\D,\;k\in\K \label{af:flow}\\
& \sum_{i\in\D} z^k_{\sigma i}\le 1
      && k\in\K \label{af:one}\\
& \sum_{i\in\D} d_i\!\sum_{a\in\delta^-(i)} z^k_a \le
      W_k\sum_{i\in\D} z^k_{\sigma i}
      && k\in\K \label{af:work}\\
& \sum_{i\in\D} e_i\, z^k_{i\omega}-\sum_{i\in\D} s_i\, z^k_{\sigma i}
      \le T_k\sum_{i\in\D} z^k_{\sigma i}
      && k\in\K \label{af:spread}\\
& z^k_a\in\{0,1\} && a\in\bar\A,\;k\in\K. \nonumber
\end{align}
Constraint~\eqref{af:cover} covers each duty exactly once,
\eqref{af:flow}--\eqref{af:one} make each employee follow at most one
$\sigma$--$\omega$ path, and \eqref{af:work}--\eqref{af:spread} enforce the
union rules. Because $k$ has exactly one first and one last duty when active,
the spread can be written \emph{without big-$M$ constants}: the two sums
in~\eqref{af:spread} pick out $e_{\text{last}}$ and $s_{\text{first}}$.
Multiplying the right-hand sides by the activation
$y_k:=\sum_i z^k_{\sigma i}$ adds nothing for integer solutions (an inactive
employee has zero left-hand side) but tightens the LP relaxation.

\begin{proposition}\label{prop:af}
For any DAG $G$, the integer solutions of (AF) correspond exactly to the
partitions of $\D$ into feasible shifts with a type assigned to each, and the
objective~\eqref{af:obj} equals the cost of the partition.
\end{proposition}
\begin{proof}
Fix $k$ with $y_k=1$. By~\eqref{af:cover}, every duty has in-flow at most $1$
under $k$. By~\eqref{af:flow} and~\eqref{af:one}, the arcs used by $k$ form a
flow of value $1$ from $\sigma$ to $\omega$. Since $G$ is acyclic, this flow
contains no cycle and is a single $\sigma$--$\omega$ path, i.e.\ a path
$p$ in $G$, which gives (S1). The left-hand side of~\eqref{af:work} is the
driving time of $p$ (S2), and that of~\eqref{af:spread} is
$e_{\text{last}}-s_{\text{first}}$ (S3). Inactive employees use no arc.
By~\eqref{af:cover} the paths partition $\D$. Conversely, every partition
into feasible typed shifts gives a feasible $z$ when the shifts are mapped
to distinct employees of the right type (the pool is large enough by the
choice of $U_\tau$).
\end{proof}

\begin{remark}
The employee index cannot simply be dropped by aggregating the flow per type:
a single integer flow decomposes into paths, but the resource
limits~\eqref{af:work}--\eqref{af:spread} can then only be imposed on the
total, not on each path. The type-aggregated flow is a relaxation of (AF); we
exploit this in Proposition~\ref{prop:lp-af}.
\end{remark}

\subsubsection{Conflict formulation (CF) --- transitive graphs}

When $G$ is transitive (assumption~\ref{ass:tri}), the order of the duties in
a shift need not be modelled: a set of duties can be sequenced if and only if
every pair of them is compatible in one direction or the other. The spread of
a set is attained by a single pair (the first start and the last end), so it
can be checked pairwise as well. For each type $\tau$ define the
\emph{conflict set}
\[
  C_\tau=\bigl\{\{i,j\}\subseteq\D:\;
     (i,j)\notin\A \text{ and } (j,i)\notin\A\bigr\}
  \;\cup\;
  \bigl\{\{i,j\}\subseteq\D:\;
     \max(e_i,e_j)-\min(s_i,s_j)>T_\tau\bigr\}.
\]
With binary variables $y_k$ (employee $k$ is used) and $x_{ik}$ (duty $i$ is
assigned to $k$):
\begin{align}
\text{(CF)}\quad
\min\;& \sum_{k\in\K} c_k\, y_k \label{cf:obj}\\
\text{s.t.}\;
& \sum_{k\in\K} x_{ik}=1 && i\in\D \label{cf:cover}\\
& x_{ik}\le y_k && i\in\D,\;k\in\K \label{cf:link}\\
& x_{ik}+x_{jk}\le y_k && \{i,j\}\in C_{\tau(k)},\;k\in\K \label{cf:conflict}\\
& \sum_{i\in\D} d_i\,x_{ik}\le W_k\,y_k && k\in\K \label{cf:work}\\
& x_{ik},\,y_k\in\{0,1\} && i\in\D,\;k\in\K. \nonumber
\end{align}
A duty with $d_i>W_k$ is automatically excluded from employee $k$
by~\eqref{cf:work}.

\begin{proposition}\label{prop:cf}
Let $G$ be transitive. A set $S\subseteq\D$ is (the duty set of) a shift
feasible for type $\tau$ if and only if no pair of $S$ lies in $C_\tau$ and
$\sum_{i\in S}d_i\le W_\tau$. Hence (CF) is a correct model of the problem.
\end{proposition}
\begin{proof}
($\Leftarrow$) Any two duties of $S$ are joined by an arc of $\A$, and since
$G$ is acyclic and transitive, $\A$ restricted to $S$ is a strict total order.
Listing $S$ in that order gives a sequence in which consecutive duties are
joined by arcs, i.e.\ a path (S1). (S2) holds by assumption. For (S3), let
$a=\arg\min_{i\in S}s_i$ and $b=\arg\max_{i\in S}e_i$. If $a=b$ the spread is
$d_a\le W_\tau\le T_\tau$; otherwise $\{a,b\}\notin C_\tau$ gives
$e_b-s_a\le T_\tau$.
($\Rightarrow$) In a feasible shift, two duties $i_\ell,i_m$ with $\ell<m$ are
joined by a path, hence by an arc by transitivity; each pair spans a sub-interval
of the shift, so its spread is at most $T_\tau$. (S2) is explicit.
\end{proof}

Without transitivity (CF) is only a relaxation: three pairwise-compatible
duties may admit no common order. In that case we use (AF).

\paragraph{Big-$M$ variant (CF-M).}
A more ``textbook'' way to model the spread replaces the second part of
$C_\tau$ by two continuous variables per employee, the start $S_k$ and end
$E_k$ of its shift:
\begin{equation}\label{cfm}
  S_k\le s_i+H(1-x_{ik}),\quad
  E_k\ge e_i\,x_{ik}\quad(i\in\D),\qquad
  E_k-S_k\le T_k\,y_k,\qquad S_k,E_k\ge 0 .
\end{equation}
This saves up to $O(n^2|\K|)$ conflict constraints, but its LP relaxation is
much weaker (with $x_{ik}$ fractional, the big-$M$ terms are nearly
inactive). We keep (CF-M) as a reference point to measure how much the exact
pairwise spread conflicts help.

\subsection{Extended formulation: set partitioning}
\label{sec:sp}

Let $\Omega=\Omega_{\mathrm F}\cup\Omega_{\mathrm H}$, where each shift
$p\in\Omega$ appears once with cost $c_p=0.6$ if $p\in\Omega_{\mathrm H}$ and
$c_p=1$ otherwise (by the remark in Section~\ref{sec:assumptions}). Let
$a_{ip}=1$ if duty $i$ belongs to $p$. With $\lambda_p=1$ if shift $p$ is
used:
\begin{align}
\text{(SP)}\quad
\min\;\sum_{p\in\Omega}c_p\lambda_p
\quad\text{s.t.}\quad
\sum_{p\in\Omega}a_{ip}\lambda_p=1\;\; (i\in\D),\qquad
\lambda_p\in\{0,1\}\;\;(p\in\Omega).
\label{sp}
\end{align}

\paragraph{Properties.}
(SP) is the Dantzig--Wolfe reformulation of (AF) (or (CF)) in which the
per-employee constraints \eqref{af:flow}--\eqref{af:spread} are convexified
and identical employees are aggregated. The symmetry between employees
disappears, and the LP bound satisfies
$z_{\mathrm{LP}}(\mathrm{SP})\ge z_{\mathrm{LP}}(\mathrm{AF})$, because the LP
of (SP) optimises over the convex hull of each employee's feasible shifts
rather than over its linear relaxation~\cite{desaulniers2005column}. When $G$
is transitive, removing a duty from a feasible shift leaves a feasible shift,
so ``$=1$'' can be relaxed to ``$\ge 1$'' (set covering) without changing the
optimal value. This makes the duals sign-constrained ($\pi_i\ge 0$) and column
generation more stable.

\paragraph{Column generation (outline, developed in Section~\ref{sec:improvements}).}
$|\Omega|$ is exponential in $n$, so the LP relaxation of (SP) is solved by
column generation. Given the duals $\pi$ of a restricted master problem, the
pricing problem asks, for each $\tau$, for a shift of minimum reduced cost
\[
  \bar c_\tau=\min_{p\in\Omega_\tau}\Bigl(c_\tau-\sum_{i\in p}\pi_i\Bigr).
\]
This is a \emph{resource-constrained shortest path problem} (RCSPP) on the DAG
$G$, with the two resources working time and spread. Paths in a DAG are
elementary, so no elementarity constraints are needed. Labels
$(\text{reduced cost},\;\text{working time},\;s_{\text{first}})$ are extended
in topological order. A label dominates another at the same node if it is no
worse in all three components (larger $s_{\text{first}}$ is better), and
labels exceeding $W_\tau$ or $T_\tau$ are discarded. With integer times the
labelling algorithm is pseudo-polynomial. An integer solution can then be
obtained by price-and-branch (solving (SP) restricted to the generated
columns as a MILP) or by branch-and-price.

\subsection{Model size and structural properties}
\label{sec:properties}

\begin{table}[ht]
\centering
\caption{Size of the formulations ($n=|\D|$, $m=|\A|$, $\K$ the employee
pool, $C=\max_\tau|C_\tau|\le n(n-1)/2$).}
\label{tab:size}
\begin{tabular}{lll}
\toprule
Model & binary variables & constraints \\
\midrule
(AF)   & $|\K|(m+2n)$ & $n+|\K|(n+3)$ \\
(CF)   & $|\K|(n+1)$  & $n+|\K|(n+C+1)$ \\
(CF-M) & $|\K|(n+1)$ (+$2|\K|$ cont.) & $n+|\K|(3n+C'+2)$, $C'\le C$ \\
(SP)   & $|\Omega|$ (exponential) & $n$ \\
\bottomrule
\end{tabular}
\end{table}

With $|\K|=\Theta(n)$ and dense graphs ($m=\Theta(n^2)$), both (AF) and (CF)
have $\Theta(n^3)$ variables or constraints. This is the main reason we expect
the compact models to stop scaling at moderate $n$.

\paragraph{Symmetry.}
Employees of the same type are interchangeable. Every solution with $f$
full-time and $h$ half-time employees appears $\binom{U_{\mathrm F}}{f}f!\,
\binom{U_{\mathrm H}}{h}h!$ times in (AF) and (CF), so branch-and-bound may
explore many equivalent subtrees. (SP) is free of this symmetry.

\paragraph{Strength of the LP relaxations.}
Let $\rho(G)$ be the minimum number of vertex-disjoint paths covering all
nodes of $G$. Every employee performs one path, so every feasible solution
costs at least $0.6\,\rho(G)$.

\begin{proposition}\label{prop:lp-af}
$z_{\mathrm{LP}}(\mathrm{AF})\ge 0.6\,\rho(G)$.
\end{proposition}
\begin{proof}
Take a feasible LP solution $z$ and set $f_a=\sum_k z^k_a$ for $a\in\bar\A$.
Summing~\eqref{af:flow} over $k$ and using~\eqref{af:cover}, $f$ is a flow in
$\bar G$ with in-flow exactly $1$ at every duty. Minimising its value
$\sum_i f_{\sigma i}$ over all such flows is the LP of the minimum path cover
problem on a DAG. After node splitting this is a network-flow problem whose
constraint matrix is totally unimodular, so its optimum is the integer value
$\rho(G)$~\cite{fulkerson1956}. Since $c_k\ge 0.6$,
$\sum_k c_k\sum_i z^k_{\sigma i}\ge 0.6\sum_i f_{\sigma i}\ge 0.6\,\rho(G)$.
\end{proof}

When $G$ is transitive, $\rho(G)$ equals the largest number of pairwise
incompatible duties, e.g.\ the maximum number of duties running
simultaneously (Dilworth's theorem~\cite{dilworth1950}). The LP of (CF) has
no such guarantee:

\begin{proposition}\label{prop:lp-cf}
There are instances with $z_{\mathrm{LP}}(\mathrm{CF})\le 2$ and optimal value
at least $0.6\,q$, for any $q$.
\end{proposition}
\begin{proof}
Take $q$ pairwise overlapping duties with $\sum_i d_i\le 2W$. Any integer
solution needs $q$ employees. In the LP of (CF), take two full-time employees
with $y_1=y_2=1$ and $x_{i1}=x_{i2}=1/2$ for every $i$:
\eqref{cf:cover}--\eqref{cf:conflict} hold, \eqref{cf:work} reads
$\sum_i d_i/2\le W$, and the cost is $2$.
\end{proof}

So (CF) is smaller but has a weak bound when many duties overlap, which is
exactly the peak-hour situation. This motivates the clique inequalities
$\sum_{i\in Q}x_{ik}\le y_k$ for sets $Q$ of pairwise conflicting duties,
which we will test in Section~\ref{sec:improvements}. For (AF) the weakness is
elsewhere: the resource constraints~\eqref{af:work}--\eqref{af:spread} are only
satisfied on average by fractional paths. (SP) removes both weaknesses.

\paragraph{Complexity.}
Without the union limits ($W,T=\infty$), half-time employees are always
preferable and the problem reduces to minimum path cover on a DAG, solvable
in polynomial time by bipartite matching~\cite{fulkerson1956}, with optimal
cost $0.6\,\rho(G)$. The limits are what make it hard:

\begin{proposition}\label{prop:np}
The problem is strongly NP-hard, even when $G$ is complete and transitive and
the spread limit is never binding.
\end{proposition}
\begin{proof}
Reduction from \textsc{3-Partition}~\cite{garey1979}: given $3m$ integers
$a_i$ with $B/4<a_i<B/2$ and $\sum_i a_i=mB$, decide whether they can be split
into $m$ triples each summing to $B$. Build duties of lengths $d_i=a_i$
placed back to back at a single location ($t_{ij}=0$), so that every pair is
compatible in time order, and set $W=B$, $T=4mB$ (so $T/2\ge H$ and the spread
never binds). We claim the optimum is $\le m$ if and only if the
\textsc{3-Partition} instance is a yes-instance. If it is, $m$ full-time
employees cover the triples at cost $m$. Conversely, let a solution use $f$
full-time and $h$ half-time employees with $f+0.6h\le m$. The total capacity
must cover the total work: $fB+hB/2\ge mB$, i.e.\ $f+0.5h\ge m$. Subtracting
gives $0.1h\le 0$, so $h=0$ and $f=m$. Then $m$ shifts of capacity $B$ hold a
total of exactly $mB$, so each sums to exactly $B$, and since
$B/4<a_i<B/2$ each holds exactly three duties.
\end{proof}

Since the pricing problem (RCSPP) is itself only weakly NP-hard and in
practice cheap on DAGs, the extended formulation moves the difficulty into a
tractable subproblem. This is the main argument for the column-generation
approach of Section~\ref{sec:improvements}.
